\section{Dynamic terminal contractions: algebra, native replay and zero cones}
\label{sec:dynamic-terminal-review}

Commit \texttt{49a8af358} supplies
\path{docs/incoming/unknot_dynamic_terminal_20261008.zip}. Its 51-file manifest
verifies. The report concerns repeated signed Tait-graph quotient observations
at a fixed suffix, extending the boundary geometry of
Section~\ref{sec:boundary-reuse}. It supplies a quadratic field-update method
for a given terminal merge, including singular intermediate matrices. It
does not supply a complete unknot criterion or bound how many observations
a complete recognition procedure must generate. We reproduced its algebra
checks, ran the previously unexecuted native adapter, and implemented its
proposed zero-descendant-cone optimization as a research primitive.

\subsection{Keep singular interior directions at each prime}

Let a signed graph have $v$ vertices and $b$ terminals, with terminal zero
grounded. Write its grounded Laplacian as
\[
 A=\begin{pmatrix}D&B\\E&F\end{pmatrix},
\]
where $D$ indexes the interior and there are $t=b-1$ non-root terminal
coordinates. Over a chosen prime field, find invertible left and right
matrices $P,Q$ such that $PDQ=\operatorname{diag}(I_\rho,0_r)$. Separate
the active and null rows of $PB$ and columns of $EQ$, and eliminate the
active identity block. The surviving matrix is
\[
 K=\begin{pmatrix}0_r&C\\G&S\end{pmatrix},
 \qquad S=F-(EQ)_{\rm active}(PB)_{\rm active}.
\]
The determinant multiplier is
$\gamma=(\det P\det Q)^{-1}$. For a partition with $q$ non-root blocks
and terminal incidence matrix $R$, its quotient cofactor is
\[
 \gamma\det\begin{pmatrix}0_r&CR\\R^T G&R^T S R\end{pmatrix}.
\]
In particular, $r>q$ forces zero. Different primes can have different $r$;
a singular interior is retained rather than discarded as a bad prime.
If $r>t$, all terminal partitions are zero at that prime. Otherwise
$\dim K=r+t\le2(b-1)$. Separate left and right elimination need not preserve
symmetry, so $G$ cannot silently be replaced by $C^T$.

\subsection{A fixed matrix and a rank-two merge}

Give each terminal block a retained anchor, with root block anchored at
zero. Let $e_i$ be the terminal unit column of $K$ and put $e_0=0$.
Keep the null-direction rows. Replace a non-root anchor row by the sum of
the \emph{original} terminal rows of $K$ in its block, and replace each
non-anchor row $i$ in anchor block $a$ by $e_i^T-e_a^T$. Denote the resulting
fixed-size matrix by $M_\pi$.

Its determinant equals the smaller quotient border above, with no extra
partition sign. Indeed set $x_i=z_a+z_i$ at non-anchor members of a non-root
block, and keep its anchor coordinate $x_a=z_a$. This column change is
$I+N$ with $N^2=0$ and determinant one. Constraint rows become unit rows.
Clear their columns, then reorder rows and columns by the same permutation;
the result is the small border direct-summed with an identity. The two
permutation signs cancel. The identity works over any commutative ring.

Now merge source block $B$ with anchor $s\ne0$ into target anchor $a$.
Define $e=e_s-e_a$, $\chi_B=\sum_{i\in B}e_i$, and let $w_B^T$ be the sum
of the original base rows indexed by $B$. Direct row comparison gives
\[
 M_{\pi'}-M_\pi=-e w_B^T+\chi_B e^T.
\]
The target anchor gains $w_B^T$; the old source anchor becomes a constraint;
every other source member changes its constraint anchor from $s$ to $a$.
The root case follows from $e_0=0$. Using current constraint rows to compute
$w_B$ would be wrong. Anchors need not be minimum labels, but a certificate
must retain their actual identities to identify the represented matrix.

\subsection{Rank updates without a nonsingularity assumption}

Maintain $L M U=\operatorname{diag}(I_k,0)$ and
$\eta\det L\det U=1$. A rank-one update transforms to $xy^T$, where
$x=Lu$ and $y^T=v^TU$. If $x$ has a nonzero null coordinate, pivot that
coordinate and use its original zero row to clear the other entries of
$x$. Clear the active part of the new row using active unit rows. A remaining
null component of $y$ adds a pivot; otherwise rank stays unchanged. The
case with a null component only in $y$ follows by transposition.

When both tails vanish, set $\alpha=1+y^Tx$. For $\alpha\ne0$ the inverse
rank-one correction restores the active identity, and $\eta$ is multiplied
by $\alpha$. If $\alpha=0$, conjugate within the active block to send $x$
to its final coordinate; its diagonal entry becomes zero and elimination
removes one pivot. All branches use $O(b^2)$ field operations and never
invert zero. The determinant is $\eta$ at full rank and zero otherwise.
Applying this routine twice realizes a terminal merge even if the first
update temporarily changes rank.

Persistent cursors clone the two normalizing matrices before either update.
The parent therefore survives an interrupted child construction. This
transactional property belongs to the cursor API; the low-level mutable
rank-update object does not promise it.

\subsection{What the local complexity bound includes}

For a supplied merge tree with $J$ edges and $Q$ observations, preparation
and evaluation take $O(v^3+Jb^2+Q)$ field and indexed operations per prime.
The cubic term includes interior preparation and the root normalization.
Depth is at most $b-1$, so depth-first traversal retains $O(v^2+b^3)$ field
elements, plus the explicit tree and results. Query generation is separate.
A static competitor should eliminate the smaller $(r+q)$ border, not the
original graph at every query. Heavy coarsening can make that static border
much cheaper than retaining a large dynamic matrix.

Signed reconstruction uses the graph-derived bound
$B=\prod_{e\text{ non-loop}}(1+|w_e|)$: every quotient spanning tree uses
a subset of the original edges. Select distinct primes with product greater
than $2B$ and reconstruct the unique signed integer in $[-B,B]$. Independent
choices of signs modulo each prime do not prove an integer is a unit:
$41$ is $-1,+1,-1$ modulo $3,5,7$, although their product exceeds $2(41+1)$.
For $H=\lceil\log_2(B+1)\rceil$, the report gives elementary polynomial
bounds on the number and size of trial-generated primes, so scalar bit
costs do not invalidate the stated explicit-graph polynomial bound.

The delivered exact tree walker reconstructs at every visited node, including
unobserved intermediate nodes. For that implementation, signed reconstruction
must be charged at $J+1$ nodes, not merely $Q$ requested observations. Our
per-prime driver reconstructs only requested outputs. This distinction does
not change the per-prime theorem, but matters for its exact-integer wrapper.
Neither implementation includes independent endpoint verification in the
quadratic update bound. Direct verification of the normal forms is cubic;
the delivered endpoint schema also repeats interior certificates.

All partitions can share a merge tree with Bell-number many nodes. This
removes repeated intermediate work but leaves
$\exp(\Theta(b\log(b+2)))$ possible partitions. Thus a suitable polylogarithmic
boundary bound can make this observer quasi-polynomial; it does not supply
such a bound for general diagrams or control the other recognition work.

\subsection{Implemented zero cones and independent native validation}

The maintained research helper
\path{../fast/determinant_research/terminal_plan.py} evaluates a supplied
prime-specific merge plan. Once $q<r$, every descendant has still fewer
blocks, so its residue is zero. The helper skips matrix cloning, updates
and algebraic traversal of that entire cone. Before pruning, it validates
\emph{every} tree edge with a depth-first live-anchor set, undoing anchor
deletions on return. An impossible merge hidden below a zero node is still
rejected. This costs $O(J+Q+b)$ indexed work even when all algebra vanishes.
The remaining algebra uses $O(v^3+J_{\rm live}b^2)$ field work, where
$J_{\rm live}$ counts actually updated edges. Requested callbacks propagate
through validation and the delivered cursor operations; unrequested work
statistics are not collected. Prime preparation and the full exact adapter
remain research interfaces without a complete production deadline contract.

The 39 delivered unit tests pass. Its deterministic audit reproduces 19,200
rank-one updates, 600 graph kernels with 193 singular interiors, 1,772 graph
quotient comparisons, all 877 partitions of seven terminals, and the three
exported certificates. Our new traversal additionally agrees with direct
integer quotient determinants on 15,840 values across 600 randomized field
plans, with nonminimum anchors and branching. Invalid merges beneath a
zero cone are rejected, and cancellation propagates.

The native driver \path{data/audit_dynamic_terminal.py} recovers the immutable
archive, rechecks its manifest and exports actual maintained scanner query
streams. On 78 diagrams and 582 stages, all 4,503 queries agree in value
\emph{and metadata} among the maintained rational observer, delivered static
and dynamic modular adapters, and independently rebuilt completed diagrams.
There are 1,317 disconnected projection queries. Connected observations also
agree with the new per-prime traversal and signed reconstruction. All 24
selected nontrivial native endpoint certificates pass independent algebraic
replay; that graph-level check is supplemented by the separate geometry
comparison, not used as a substitute for it.

Across 1,819 prepared prime kernels, 132 interiors are singular and 24 are
universally zero. The connected query plans contain 4,535 merge edges.
Beyond universal zeros, the new pruning reduces field-edge updates from
16,862 to 16,785: only 77 are avoided. This is a real activation on native
inputs, but a small reduction, not evidence for a broadly effective shortcut.
These stages have at most ten graph vertices, ten terminals and 200 queries;
their field-interior nullities are at most two. They therefore test actual
integration but do not probe the large-boundary asymptotic regime of the
report's graph microbenchmarks.

\subsection{Controlled native-stage measurements}

\begin{center}\small
\begin{tabular}{@{}lrrrrr@{}}
\toprule
Observer & All stages & Widest & Most queries & Singular & Ratio \\
\midrule
Rational & 351.585 & 1.366 & 10.829 & 1.007 & 1.00 \\
Identical control & 360.448 & 1.323 & 9.815 & 0.999 & 0.97 \\
Integer quotient & 293.190 & 1.299 & 7.278 & 0.663 & 1.24 \\
Static modular & 506.603 & 1.770 & 12.928 & 1.731 & 0.69 \\
Cached static modular & 514.974 & 1.969 & 13.275 & 1.635 & 0.70 \\
Delivered dynamic & 1276.307 & 7.055 & 35.974 & 2.638 & 0.28 \\
Pruned field traversal & 1406.933 & 6.993 & 37.377 & 2.837 & 0.27 \\
\bottomrule
\end{tabular}
\end{center}
Median milliseconds per fresh replay; the final column is the median paired rational/arm ratio for all 582 stages. These are observer timings.


The final benchmark includes fresh PD validation, boundary geometry, kernel
preparation, matching validation, plan compilation and exact reconstruction.
It excludes archive import, output hashing and the separate independent
audits. The rational reference and its identical control retain their
normalized partition cache. The integer arm constructs full integer quotient
cofactors and uses fraction-free elimination with the same cache. We measure
both delivered static modular evaluation and a cached version, delivered
dynamic evaluation, and the new per-prime pruned traversal. The last arm
also avoids reconstruction at unrequested tree nodes, so its timing is not
an isolated measurement of the zero-cone test.

One excluded warmup batch and five shuffled measured rounds use seed
261008116. The complete 582-stage corpus uses one replay per sample; widest,
most-query and most-singular stages use respectively 40, 10 and 50 fresh
replays per batch. All repeated results must have identical hashes including
metadata. The initial unbatched record is retained separately: its smallest
case had a 1.78 identical-control ratio, and its new traversal unnecessarily
collected work counters. The final record uses optional counters and batches
to address those concrete measurement problems. Frozen initial source and
audit snapshots permit verification of both records.

All 140 measured batches and 28 warmup batches complete with matching output
hashes. Across the full corpus the rational median is 351.6 ms, its identical
control is 360.4 ms, and the integer quotient median is 293.2 ms; the paired
rational/integer ratio is 1.245. Delivered dynamic evaluation takes 1,276.3 ms
with paired rational/dynamic ratio 0.278, roughly a 3.6-fold slowdown.
The new traversal takes 1,406.9 ms and supplies no aggregate timing gain.
Static modular evaluation is also slower, with or without a partition cache.
The final identical-control paired ratios range from 0.970 to 1.033 over
the four workloads, much closer than in the initial small-case measurement.
These results motivate testing integer/rational adaptation across larger
native stages before choosing a new policy; they do not establish a
competitive adaptive bound or justify a cutoff fitted to this small corpus.

These are boundary-observer measurements, not full-recognizer speedups.
The native adapter and zero-cone helper remain research-only. The report's
general local theorem is useful, but the native evidence does not justify
replacing maintained dispatch with an eager dynamic modular observer.
